\scanpage{21}

Donc on a les relations

\[
\left\{\begin{aligned}
\widetilde\sigma_0^2&=\widetilde\sigma_1^2=\widetilde\sigma_\infty^2=1\\
\widetilde\sigma_0\widetilde\sigma_1&=\widetilde\sigma_1\widetilde\sigma_\infty=\widetilde\sigma_\infty\widetilde\sigma_0=\rho^{\deleted{\unsure{\sim}}}\deleted{\unsure{(=\rho)}}\\
(\rho^{\deleted{\unsure{\sim}}\,3}&=1\ \text{ est conséquence des précédentes})
\end{aligned}\right. \tag*{(13)}
\]

% slanted
\margin{NB $\rho\widetilde\sigma_i\rho^{-1}=\widetilde\sigma_{\dot\rho i}$ i.e\\
$\rho\widetilde\sigma_0\rho^{-1}=\widetilde\sigma_1$\\
$\rho\widetilde\sigma_1\rho^{-1}=\widetilde\sigma_\infty$\\
$\rho\widetilde\sigma_\infty\rho^{-1}=\widetilde\sigma_0$}

et l'image inverse $\unsure{\widetilde\Sigma|\widetilde E}$ de $\varphi({\unsure{\mathfrak S_3|\mathcal S_3}}^{\deleted{\unsure{\sim}}})\subset\Gamma$ (ss-groupe d'indice 2) est produit \unsure{1/2|s/d} direct de $\mathfrak S_3$ et de $\pi$. Il faut expliciter les opérations des $\widetilde\sigma_i$ sur $\pi$ pour avoir la structure complète du groupe, On trouve que $\widetilde\sigma_{\unsure{i|0}}$ est au-dessus de $\varphi(\dot\sigma_{\unsure{i|0}}^{\deleted{\unsure{\sim}}})$ i.e. la symétrie p.r. au plan passant par $P,\overline P,\deleted{\unsure{\infty}}\,i$, qui transforme $l_0$ en $l_0^{-1}$, $\deleted{\unsure{l}}l_\infty$ en $l_{\unsure{1}}^{-1}$, $l_1$ en $l_\infty^{-1}$ :

% slanted
\margin{Ce sont \unsure{elles|celles} \unsure{se trouvent par|obtenues par} transport de structure, via les opérations de $\mathfrak S_3$ sur $(\unsure{U|V},P)$}

% slanted
\margin{\uline{Critique} La notation $\widetilde\sigma_0$ suggère \unsure{fautivt} que $\widetilde\sigma_0$ est au-dessus de $\sigma_0$, -- il vaut mieux écrire $\dot\sigma'_i=\dot\sigma_i\tau$ et noter $\widetilde\sigma'_i$ au lieu des $\widetilde\sigma_i$ $(i\in\{0,1,\infty\})$, ou mieux $\sigma'_i$ (sans $\sim$) -- pour réserver $\widetilde\sigma'_i$ à des éléments de $S\Sigma$, cf plus bas \ldots}

\[
\left\{\begin{aligned}
\widetilde\sigma_0 l_0\widetilde\sigma_0^{-1}&=l_0^{-1}&
\widetilde\sigma_0 l_1\widetilde\sigma_0^{-1}&=l_\infty^{-1}&
\widetilde\sigma_0 l_\infty\widetilde\sigma_0^{-1}&=l_1^{-1}\\
\widetilde\sigma_1 l_1\widetilde\sigma_1^{-1}&=l_1^{-1}&\ldots\\
&\ldots
\end{aligned}\right. \tag*{(14)}
\]

Les relations (13), (14), plus la relation fondamentale $l_\infty\inserted{\unsure{\sim|\checkmark}}l_1l_0=1$ (1), font une présentation de $\unsure{\widetilde G|\widetilde\Gamma}$, dont on peut si on veut éliminer $\rho$, en omettant d'écrire la relation qui le \unsure{définit|définist}\unsure{)|.}

On peut maintenant remonter $\dot\tau$ en trois manières, \unsure{à|6} \unsure{par|peu}

\[
\left\{\begin{aligned}
\tau_0&=\sigma_0\widetilde\sigma_0=\widetilde\sigma_0\sigma_0\\
\tau_1&=\sigma_1\widetilde\sigma_1=\widetilde\sigma_1\sigma_1\\
\tau_\infty&=\sigma_\infty\widetilde\sigma_\infty=\widetilde\sigma_\infty\sigma_\infty
\end{aligned}\right. \tag*{(15)}
\]

\deleted{donc}

\[
\left\{\begin{aligned}
\tau_0^2&=1\\
\tau_1^2&=1\\
\tau_\infty^2&=1
\end{aligned}\right. \tag*{(16)}
\]
% archive sheet number: 20
